\documentclass[10pt,a4paper]{amsart}
\usepackage[margin=19mm]{geometry}
\usepackage{amsmath,amssymb,mathtools}
\usepackage[scr=rsfs,cal=euler]{mathalfa}
\usepackage{xfrac}
\usepackage[unicode,hidelinks,bookmarks=false]{hyperref}
\newcommand{\ie}{i.e.\ }
\newcommand{\cf}{cf.\ }
\newcommand{\resp}{respectively\ }
\newcommand{\Subsection}{\subsection}
\providecommand{\nobreakdash}{\nobreak\hbox{-}}
\setlength{\emergencystretch}{2em}
\multlinegap0pt
\usepackage{dsfont}
\usepackage{amssymb}
\usepackage{bbm}
\usepackage{mathtools}
\newtheorem{thm}{Theorem}
\newcommand{\Prob}{\mathbb{P}}
\newcommand{\R}{\mathbb{R}} 
\DeclareMathOperator{\Var}{Var}
\title{A note on the Littlewood-Offord problem for discrete log-concave distributions}
\author{Arnaud Marsiglietti, James Melbourne}
\date{Source: arXiv:2510.25869v2 (CC BY 4.0)}
\begin{document}
\maketitle

\noindent\emph{Source and licence.} A note on the Littlewood-Offord problem for discrete log-concave distributions. Version arXiv:2510.25869v2 (CC BY 4.0), distributed by arXiv under the Creative Commons Attribution 4.0 International licence, http://creativecommons.org/licenses/by/4.0/, as recorded in the arXiv licence metadata for this exact version and shown on the abstract page https://arxiv.org/abs/2510.25869v2. Full licence notice: \texttt{licenses/arxiv-2510.25869-CC-BY-4.0.txt}.\par\smallskip

\noindent\textit{Editorial scope.} Counted unit: the article's thm main3 (source0.tex lines 276--282). The article's complete original proof of it is delivered with it (source0.tex lines 623--639, one \texttt{proof} environment of 17 lines, which the article places away from the statement). No mathematics is written, repaired or re-derived: the statement and the proof are the article's own lines, in the article's own notation, and no retained line is substituted or re-flowed.  Retained uncounted as the source-local context the proof needs: lines 243--251.  Nothing was omitted from any retained span, so the package carries no omission record.  No retained line contains a citation, so the wrapper carries no bibliography and no bibliography entry.  No retained line contains a figure environment, an image-inclusion command or drawing code, so no image is delivered and the package declares no asset dependency.  Editorial formatting only, in addition: an AMS \texttt{amsart} preamble that restates the article's own declarations and macros used by the retained lines, namely the environments \texttt{thm} on separate counters, together with the standard packages those lines need.  The retained environments print in this document as Theorem 1 in order of appearance, whereas the article prints the same items with the numbering of its own full document.  The complete native source of the article as distributed in the arXiv e-print for this exact version is bundled unchanged as \texttt{sources/188075/source0.tex}.
\begin{thm}\label{main1}

Let $X_1, \dots, X_n$, $n \geq 1$, be independent discrete log-concave random variables finitely supported. Then,
\begin{equation}\label{1}
\sup_{a \in (\R \setminus \{0\})^n} \sup_{x \in \R} \, \Prob(a_1 X_1 + \cdots + a_n X_n = x) \leq \frac{1}{\sqrt{1 + c \sum_{k=1}^n \Var(X_k)}},
\end{equation}
with $c=1$. Moreover, one may take $c = 2$ when the random variables are, in addition, symmetric about a point.

\end{thm}

\begin{thm}\label{main3}
Let $X_1, \dots, X_n$, $n \geq 1$, be independent discrete log-concave random variables finitely supported. Then,
\begin{equation*}
\sup_{a \in (\R \setminus \{0\})^n} \sup_{x \in \R} \, \Prob(a_1 X_1 + \cdots + a_n X_n \in A_{l,m}(x)) \leq \frac{l}{\sqrt{1 + c \sum_{k=1}^n \Var(X_k) + c \, \frac{l^2-1}{12}}},
\end{equation*}
with $c=1$. Moreover, one may take $c = 2$ when the random variables are, in addition, symmetric about a point.
\end{thm}

\begin{proof}[Proof of Theorem \ref{main3}]
Let $Y$ be a discrete random variable independent of $U_l$, where $U_l$ is uniform on $\{1,2,\dots, l\}$. Then,
\begin{align*}
    \mathbb{P}(Y - m U_l = x) 
        &= 
            \sum_{k=1}^l \mathbb{P}(U_l = k, Y = x + mk ) 
                \\
        &=
            \frac 1 l \sum_{k=1}^l \mathbb{P}(Y = x +mk)
                \\
        &=
            \frac 1 l \mathbb{P}( Y \in A_{l,m}(x)).
\end{align*}
Thus, this Littlewood-Offord problem for arithmetic progressions is equivalent to finding upper bounds on $M(a \cdot X - m U_l)$. When the $X_k$'s are discrete log-concave, one may thus apply Theorem \ref{main1} to obtain
$$ \sup_{x \in \R} \, \mathbb{P}(a \cdot X \in A_{l,m}(x)) = l \, M(a \cdot X - m U_l) \leq \frac{l}{\sqrt{1 + c \left( \sum_{k=1}^n \Var(X_i) + \Var(U_l) \right)}}, $$
which is the desired result since $\Var(U_l) = (l^2-1)/12$.
\end{proof}

\end{document}
